We set out to test whether the error of a feature attribution, measured against the causal effect it is commonly read as, can be predicted without causal ground truth. On a benchmark of 300 synthetic causal worlds the answer has two halves. Most of the error is structural: exact Shapley values of the true outcome function already misallocate a third of the importance mass relative to the total interventional effect, and better explainers do not change that. The part of the error that observable signals can see is the part contributed by the fitted model and the explainer; a gradient-boosted diagnostic predicts those components well, transfers to graph families it never saw, and beats every single-signal practitioner baseline for the total share error. The structural estimand gap, however, is largely invisible to the same signals, and the feature ranking error that follows from it is hardly predictable. The practical conclusion is a modest one: observable diagnostics can tell a practitioner when to distrust an explanation because the model or the explainer is at fault, but they cannot substitute for the causal knowledge that decides whether a Shapley attribution answers the interventional question at all. Code, the benchmark, the notebook and all tables are available at \url{https://github.com/stacktrace-bhuvi/datasphere}.